\begin{tikzpicture}[scale=0.777]

\draw[->] (-0.3,0) -- (8,0) node[right] {$x$} node[below right] {$(1/n)$};
\draw[->] (0,-0.3) -- (0,7) node[above] {$y$};

\draw[color=blue,dashed,thick] plot[domain=0:3,samples=100,variable=\x] ({\x},{(1+\x)^1.5}) node[above] {$(1+x)^t,\ t>1$};

\draw[color=blue,dashed,thick] plot[domain=0:6,samples=100,variable=\x] ({\x},{(1+\x)^0.5}) node[above right] {$(1+x)^t$} node[right] {$0<t<1$};

\foreach \t/\tdomain in {0.8/6.3,1.13/5,1.3/4,2/1.8}{
\draw[color=blue,dashed,thick] plot[domain=0:\tdomain,samples=100,variable=\x] ({\x},{(1+\x)^(\t)});
}

\path[color=blue,bend right,->] (8,6) edge node[midway,above,sloped] {\color{blue}$t$增加} (6,8);

\only<2->{
\draw[color=themecolor,dashed,thick] plot[domain=0:6,samples=100,variable=\x] ({\x},{1+\x}) node[below right] {$1+x$};
}

\only<3->{
\draw[color=red,dashed,thick] plot[domain=0:3.1,samples=100,variable=\x] ({\x},{1+2*\x}) node[above right] {$1+qx$};
}

\only<1-3>{
\begin{scope}[on background layer]
\draw[color=zkbackground] (0.8,2.8) circle (2.1);
\end{scope}
}

\only<4->{
\begin{scope}[on background layer]
\draw[black,dashed,fill=red!10] (0.8,2.8) circle (2.1);
\end{scope}
}

\only<1-4>{
\node at (9,4.7) {\phantom{\color{red}可以从函数图像上看到，$1+qx$ 的值的确可能在自变量 $x=0$ 的某个邻域内大于 $(1+x)^t$ 的值，$t>1$}};
}

\only<5->{
\node[draw=black,text width=0.47\textwidth,fill=cyan!30] at (9,4.7) {\color{red}可以从函数图像上看到，$1+qx$ 的值的确可能在自变量 $x=0$ 的某个邻域内大于 $(1+x)^t$ 的值，$t>1$};
}

\end{tikzpicture}
